\documentclass[a4paper]{article}
% Español (México) con XeLaTeX: fontspec para fuentes Unicode, polyglossia
% para la división silábica y los nombres del documento en la variante mexicana.
\usepackage{fontspec}
\usepackage{polyglossia}
\setdefaultlanguage[variant=mexican]{spanish}
% Los nombres chinos van como «pinyin (original)»: xeCJK da a los caracteres
% Han su propia fuente; sin ella XeLaTeX los omite con un aviso.
% FandolSong no cubre algunos caracteres de nombres (U+73FA); WenQuanYi Zen Hei sí.
\usepackage[AutoFallBack=true]{xeCJK}
\setCJKmainfont{FandolSong-Regular.otf}
\setCJKfallbackfamilyfont{\CJKrmdefault}{WenQuanYi Zen Hei}
% polyglossia no traduce los nombres de listings.
\AtBeginDocument{\providecommand{\lstlistingname}{}\renewcommand{\lstlistingname}{Listado}%
  \providecommand{\lstlistlistingname}{}\renewcommand{\lstlistlistingname}{Índice de listados}}
\usepackage{amsmath, amssymb}
\usepackage{graphicx}
\usepackage[margin=2.35cm]{geometry}
\usepackage[most]{tcolorbox}
\usepackage{booktabs}
\usepackage{tabularx}
\usepackage{longtable}
\usepackage{array}
\usepackage{float}
\usepackage{xcolor}
\usepackage{hyperref}
\usepackage{fancyhdr}
\setCJKmainfont[AutoFakeBold=2.4]{AR PL UMing CN}
\setCJKsansfont[AutoFakeBold=2.4]{Droid Sans Fallback}
\setCJKmonofont[AutoFakeBold=2.4]{Droid Sans Fallback}
\hypersetup{colorlinks=true, linkcolor=blue!55!black, urlcolor=blue!60!black}
\graphicspath{{./}{figures/}}
\setlength{\parskip}{0.55em}
\setlength{\parindent}{2em}
\setlength{\emergencystretch}{3em}
\renewcommand{\arraystretch}{1.18}
\newtcolorbox{knowledgebox}[1]{enhanced, colback=blue!5!white, colframe=blue!70!black, colbacktitle=blue!70!black, coltitle=white, fonttitle=\bfseries, title={#1}, sharp corners, breakable}
\newtcolorbox{importantbox}[1]{enhanced, colback=yellow!10!white, colframe=yellow!75!black, colbacktitle=yellow!75!black, coltitle=black, fonttitle=\bfseries, title={#1}, sharp corners, breakable}
\newtcolorbox{warningbox}[1]{enhanced, colback=red!5!white, colframe=red!70!black, colbacktitle=red!70!black, coltitle=white, fonttitle=\bfseries, title={#1}, sharp corners, breakable}
\pagestyle{fancy}
\fancyhf{}
\fancyfoot[C]{\thepage}
\fancyfoot[R]{\footnotesize\color{gray}\href{https://github.com/hqhq1025/ai-course-notes}{AI Course Notes}}
\renewcommand{\headrulewidth}{0pt}
\renewcommand{\footrulewidth}{0pt}
\newcommand{\videourl}{https://www.youtube.com/watch?v=gQgKkUsx5q0}
\newcommand{\repourl}{https://github.com/hqhq1025/ai-course-notes}

\begin{document}
\begin{titlepage}
\centering
{\Large Notas didácticas de una entrevista a fondo\par}
\vspace{1.0cm}
{\huge\bfseries The Second Half: Agents de lenguaje, mecanismos de recompensa y los límites de la inteligencia}\par
\vspace{0.5cm}
{\Large Zhang Xiaojun conversa con Yao Shunyu (姚顺雨), investigador de OpenAI\par}
\vspace{0.35cm}
{\large Elaborado a partir del video público de Zhang Xiaojun Podcast, la transcripción hecha en GPU, la estructura de capítulos y diagramas conceptuales propios\par}
\vspace{0.3cm}
{\large 2025-09-11\par}
\vspace{0.85cm}
\includegraphics[width=0.88\textwidth,height=0.42\textheight,keepaspectratio]{cover.jpg}\par
\vfill
\begin{tcolorbox}[width=0.92\textwidth, colback=black!2!white, colframe=black!55, sharp corners]
\textbf{Título del video}: 115. Entrevista de 3 horas con Yao Shunyu de OpenAI: 6 años de investigación sobre Agents, las personas y los sistemas, la frontera de lo que se devora, un mundo a la vez unipolar y plural\par
\textbf{Canal del video}: Zhang Xiaojun Podcast\par
\textbf{Duración del video}: 02:31:32\par
\textbf{Enlace del video}: \href{\videourl}{\nolinkurl{\videourl}}\par
\textbf{Notas sobre el material}: YouTube no ofrece subtítulos; el cuerpo del texto se elaboró a partir de la transcripción con faster-whisper large-v3 en CUDA, los capítulos del video, la descripción del video, la portada y figuras didácticas propias. El video con portada fija no se reutiliza en el cuerpo del texto.\par
\textbf{Página del proyecto}: \href{\repourl}{\nolinkurl{\repourl}}\par
\end{tcolorbox}
\end{titlepage}

\tableofcontents
\newpage

\section{Introducción: la segunda mitad de la IA pasa de la capacidad del modelo a la capacidad del sistema}

Esta sección establece primero las coordenadas de todo el episodio. Yao Shunyu (姚顺雨) publicó «The Second Half» en abril de 2025, donde sostiene que el hilo principal de la IA ya entró en su segunda mitad. La línea principal de la primera mitad se parece más a la capacidad del modelo: preentrenamiento, escalamiento, conversación y capacidades generales; la segunda mitad se parece más a la capacidad del sistema: Agent, entornos de tareas, mecanismos de recompensa, llamadas a herramientas, interfaces de interacción y estructura organizacional. La entrevista parte de la experiencia personal, pero su núcleo real es una pregunta: cuando el modelo es lo bastante potente, ¿cómo entra la inteligencia en el mundo y da lugar a acción, organización y nuevas fronteras?

\begin{figure}[H]
\centering
\includegraphics[width=0.92\textwidth]{second-half-map.png}
\caption{Mapa de la segunda mitad de la IA: de la capacidad del modelo hacia Agent, sistema, interacción y el panorama humano global. Diagrama conceptual propio, elaborado a partir de la conversación de 00:01:45--02:31:20.}
\end{figure}

\begin{knowledgebox}{Lectura de la figura: la segunda mitad no significa «el modelo ya no importa»}
El modelo sigue siendo la base, pero el centro de la discusión pasa de la puntuación de un modelo individual a cómo actúa el sistema: si la tarea es real, si el entorno es lo bastante amplio, si la reward puede definirse, si la interface abre nuevas oportunidades y si el mundo final será unipolar o plural.
\end{knowledgebox}

\begin{importantbox}{Tesis central del episodio}
Un Agent no es un botón de función de un modelo conversacional, sino un paradigma de sistema: debe observar, razonar, actuar y recibir retroalimentación dentro de un entorno, y ampliar continuamente la frontera de las tareas que puede completar mediante herramientas, recompensas y estructuras multiagente.
\end{importantbox}

\begin{knowledgebox}{Tabla general de conceptos clave}
\begin{tabularx}{\textwidth}{p{0.20\textwidth}p{0.34\textwidth}X}
\toprule
Concepto & Significado en este episodio & Malentendido que conviene evitar al estudiar \\
\midrule
Agent & Sistema que interactúa con un entorno e intenta optimizar un objetivo & No es tan simple como «una ventana de chat con algunas herramientas». \\
Environment & El mundo donde el Agent actúa y recibe retroalimentación & Si el entorno es demasiado pequeño, la capacidad no puede extrapolarse. \\
Reward & Señal que evalúa las consecuencias de una acción & La reward no equivale a todo el valor que las personas realmente quieren. \\
Affordance & Posibilidades de acción que el entorno ofrece a quien actúa & El code/API es la mano de la IA, no solo un formato de salida. \\
Interface & Forma en que el usuario interactúa con el sistema inteligente & La UI cambia las tareas, los datos, el negocio y la dirección de la investigación. \\
Different Bet & Una ruta de apuesta distinta a la del actor dominante & No es copiar con otra apariencia, sino cambiar las condiciones de frontera. \\
\bottomrule
\end{tabularx}
\end{knowledgebox}

\begin{knowledgebox}{Ruta de lectura del episodio}
\begin{tabularx}{\textwidth}{p{0.22\textwidth}p{0.32\textwidth}X}
\toprule
Ruta & Pregunta central & Secciones correspondientes \\
\midrule
Persona & ¿Por qué el lenguaje es una herramienta de generalización y por qué Yao Shunyu apuesta por Agent? & Lenguaje y posturas no consensuadas. \\
Sistema & ¿Cómo evolucionan juntas las líneas de tareas, entornos y métodos de Agent? & Las dos preguntas centrales y las tres oleadas de auge y declive. \\
Recompensa & ¿Cómo obtiene un Agent reward, autoexploración y estructura multi-agent? & Reward, code affordance, generalización. \\
Frontera & ¿Cómo determinan Chatbot, interface y Super App las nuevas oportunidades? & La absorción de fronteras y las formas de interacción. \\
Panorama global & ¿Cómo coexisten lo unipolar y lo plural, OpenAI y las oportunidades de emprendimiento? & El panorama humano global y different bet. \\
\bottomrule
\end{tabularx}
\end{knowledgebox}

\subsection{Resumen de la sección}

EP115 es una entrevista sobre la teoría de Agent, no solo una entrevista a un investigador de OpenAI. Reúne en un mismo mapa el lenguaje, las personas, las tareas, los entornos, las recompensas, la interacción y la competencia entre plataformas, por lo que sirve como nota central para entender la segunda mitad de los Agent.
\section{Las personas: el lenguaje es una herramienta inventada para generalizar}

Esta sección examina primero la parte de «las personas». Una afirmación central de Yao Shunyu (姚顺雨) es que el lenguaje es una herramienta que los seres humanos inventaron para lograr la generalización, y que esto es más fundamental que cualquier otra cosa. Esta idea explica por qué él partió del lenguaje para trabajar en Agent: el lenguaje no es solo un medio de comunicación, sino también la herramienta con la que los seres humanos comprimen la experiencia, transmiten reglas, describen objetivos, organizan planes y transfieren lo aprendido de un escenario a otro.

\begin{figure}[H]
\centering
\includegraphics[width=0.94\textwidth]{language-generalization.png}
\caption{El lenguaje como herramienta de generalización: los seres humanos comprimen la experiencia con el lenguaje, y el Agent usa el lenguaje para conectar las tareas con el mundo. Diagrama conceptual de elaboración propia, basado en la conversación de 00:02:03--00:07:50.}
\end{figure}

\begin{knowledgebox}{Lectura de la figura: por qué el lenguaje es adecuado como capa intermedia del Agent}
El lenguaje convierte las tareas, las reglas, las restricciones, las descripciones de herramientas y la retroalimentación en símbolos que pueden combinarse. Mediante el lenguaje, el modelo puede comprender los objetivos humanos y también invocar código, API, documentos y otros Agent. Por ello, el lenguaje no es todo el Agent, pero sí es una interfaz general muy potente.
\end{knowledgebox}

\subsection{Una postura fuera del consenso: apostar desde el principio por los Agent de lenguaje}

La sección anterior explicó por qué el lenguaje es una herramienta de generalización; esta examina cómo ese juicio se convirtió en una apuesta de investigación. Yao Shunyu dice que siempre sostuvo una postura fuera del consenso: quería trabajar en Agent y, como primer paso, construirlos sobre modelos de lenguaje. En ese momento, la corriente dominante no necesariamente consideraba que los modelos de lenguaje fueran el mejor camino hacia los Agent; pero después la serie GPT mostró que los modelos de lenguaje tienen un gran potencial de generalización y de servir como interfaz con herramientas, y esa postura fuera del consenso pasó a ser la línea principal de la frontera de investigación.

\begin{importantbox}{El sentido de una Different Bet}
Si solo se copia la línea principal del líder ya establecido, es muy difícil superarlo. Las nuevas oportunidades suelen surgir de una different bet: tareas distintas, entornos distintos, interface distinta, formas de producto distintas. En sus inicios, los Agent de lenguaje fueron justamente una apuesta de este tipo.
\end{importantbox}

\begin{knowledgebox}{Por qué los primeros Agent de lenguaje estaban fuera del consenso}
Los primeros modelos de lenguaje parecían más sistemas de texto que agentes inteligentes capaces de actuar. Para verlos como Agent había que creer tres cosas a la vez: que el lenguaje puede expresar tareas, que el modelo puede generalizar de una tarea a otra y que las herramientas externas pueden convertir los planes expresados en lenguaje en acciones. En ese momento, ninguna de las tres era un consenso evidente.
\end{knowledgebox}

\begin{knowledgebox}{Lectura didáctica de la trayectoria de investigación de Yao Shunyu}
\begin{tabularx}{\textwidth}{p{0.20\textwidth}p{0.34\textwidth}X}
\toprule
Etapa & Problema de interés & Relación con el Agent \\
\midrule
Inicio del doctorado & Usar modelos de lenguaje para juegos y tareas sencillas & Demostrar que el lenguaje puede contener decisiones y estados. \\
Entornos de tareas & Buscar benchmark más realistas y valiosos & Lograr que el Agent no se limite a resolver ejercicios pequeños. \\
Herramientas/código & InterCode, entornos de código, affordance del mundo digital & Darle al Agent manos capaces de ejecutar. \\
Segunda mitad & reward, multi-agent, interface, evolución de sistemas & Pasar de la capacidad de un solo modelo a la capacidad del sistema. \\
\bottomrule
\end{tabularx}
\end{knowledgebox}

\subsection{Resumen de la sección}

La sección sobre «las personas» muestra que la investigación sobre Agent no surgió de la nada a partir de las capacidades de los modelos, sino de una comprensión del lenguaje, de la generalización humana y de la expresión de tareas. El lenguaje es importante porque conecta los objetivos humanos con la acción de las máquinas.
\section{Sistemas: tareas, entornos y métodos simples y generales}

La sección anterior explicó por qué el lenguaje es la puerta de entrada; esta sección aborda el núcleo del Agent como sistema. Yao Shunyu (姚顺雨) resume su investigación en dos preguntas: primero, cómo construir tareas y entornos valiosos y más vinculados con el mundo real; segundo, cómo construir métodos simples pero generales. La primera es la línea de las tareas; la segunda, la línea de los métodos. Muchas discusiones solo miran la línea de los métodos y descuidan la de las tareas, pero las capacidades de un Agent suelen estar definidas por el entorno.

\begin{figure}[H]
\centering
\includegraphics[width=0.96\textwidth]{agent-two-core-questions.png}
\caption{Los dos ejes de la investigación sobre Agents: tareas y entornos valiosos + métodos simples y generales. Diagrama conceptual propio, elaborado a partir de la conversación entre 00:17:50 y 00:35:00.}
\end{figure}

\begin{knowledgebox}{Lectura de la figura: tareas y métodos deben evolucionar juntos}
Si la tarea es demasiado simple, el método parece muy potente pero carece de relevancia práctica; si la tarea es lo bastante realista pero el método es demasiado complejo, resulta difícil generalizarlo. La investigación sobre Agents debe diseñar a la vez el entorno, la recompensa, las herramientas y los métodos generales.
\end{knowledgebox}

\subsection{Las tres olas de auge y declive de los Agents: la línea de los métodos y la línea de las tareas}

La sección anterior afirmó que tareas y métodos deben evolucionar juntos; esta sección sitúa ese juicio en la historia de los Agents. Agent es un concepto antiguo: cualquier sistema capaz de tomar decisiones por sí mismo, interactuar con un entorno y optimizar una recompensa puede llamarse Agent. La entrevista nos recuerda que la historia de los Agents no es una línea recta, sino una sucesión de auges y declives: los Agents de la IA clásica y del RL, los entornos de aprendizaje por refuerzo y, después, los LLM Agents. Cada ola no fue solo un cambio de métodos; también dependió de que las tareas y los entornos fueran lo bastante adecuados.

\begin{figure}[H]
\centering
\includegraphics[width=0.96\textwidth]{agent-history-waves.png}
\caption{Las tres olas de auge y declive de los Agents: la línea de los métodos y la de las tareas se complementan; no basta con mirar los algoritmos. Diagrama conceptual propio, elaborado a partir de la conversación entre 00:17:50 y 00:29:00.}
\end{figure}

\begin{knowledgebox}{Términos clave: componentes básicos de un sistema de Agent}
\begin{tabularx}{\textwidth}{p{0.20\textwidth}p{0.32\textwidth}X}
\toprule
Componente & Significado & Por qué importa \\
\midrule
Environment & El mundo que el Agent observa y en el que actúa & Determina si la tarea es realista y si la retroalimentación es válida. \\
Policy & La estrategia que elige acciones según las observaciones & Puede implementarse con un LLM, código, un planificador o un sistema de varios modelos. \\
Reward & La señal que evalúa el resultado de las acciones & Determina qué aprende y qué explora el Agent. \\
Memory & Estado de largo plazo y registro de experiencias & Sostiene el trabajo entre tareas, entre sesiones y la personalización. \\
Tools & Capacidades externas como API, código, navegador y archivos & Permiten que el Agent pase del lenguaje a un mundo ejecutable. \\
\bottomrule
\end{tabularx}
\end{knowledgebox}

\begin{knowledgebox}{Línea de los métodos vs línea de las tareas}
\begin{tabularx}{\textwidth}{p{0.22\textwidth}p{0.34\textwidth}X}
\toprule
Línea & En qué se enfoca & Qué suele pasar por alto \\
\midrule
Línea de los métodos & Algoritmos, modelos, planificación, técnicas de entrenamiento & Si la tarea es lo bastante realista y si el entorno puede dar retroalimentación. \\
Línea de las tareas & benchmark, entornos, herramientas, datos y evaluación & Si el método es simple y general, y si puede transferirse a tareas nuevas. \\
Línea del sistema & Cómo métodos y tareas forman un ciclo cerrado & Costo de cómputo, recuperación ante fallos, memoria de largo plazo y seguridad. \\
\bottomrule
\end{tabularx}
\end{knowledgebox}

Esta tabla también explica por qué muchas demos de Agents parecen impresionantes y, sin embargo, difícilmente se convierten en productos estables. Una demo suele mostrar solo la línea de los métodos: el modelo planifica, invoca herramientas y escribe algunos pasos de acción; un producto real también debe resolver la línea de las tareas: si el entorno es reproducible, si la retroalimentación puede evaluarse de forma inequívoca, si los fallos son recuperables y si el usuario está dispuesto a confiarle sus objetivos reales. Cuando Yao Shunyu insiste en las tareas y los entornos, recuerda precisamente que la investigación sobre Agents no puede limitarse a «aparentar que sabe hacer las cosas».

\begin{importantbox}{Criterios de diseño de un benchmark de Agents}
Un buen benchmark de Agents debe cumplir al menos tres condiciones: tareas lo bastante realistas, retroalimentación lo bastante clara y un entorno lo bastante abierto. Las tareas demasiado simples no permiten medir la generalización; los entornos demasiado cerrados no permiten entrenar acciones reales, y las tareas sin retroalimentación confiable no generan una señal de aprendizaje.
\end{importantbox}

\subsection{Code es la affordance clave}

Hasta aquí se discutieron los componentes de un Agent; esta sección aborda la «interfaz de acción». En la entrevista aparece un juicio muy importante: el code es como la mano humana, la affordance más importante de la IA. Affordance designa las posibilidades de acción que el entorno ofrece a quien actúa. Para la IA, el código, las API, la línea de comandos, el navegador y el sistema de archivos permiten que el modelo no solo hable, sino que modifique el mundo digital.

\begin{figure}[H]
\centering
\includegraphics[width=0.96\textwidth]{code-affordance.png}
\caption{Code es la Affordance clave: el código permite que la IA pase del lenguaje a la acción en el mundo digital. Diagrama conceptual propio, elaborado a partir de la conversación entre 00:29:49 y 00:33:00.}
\end{figure}

\begin{importantbox}{Por qué el code es la mano}
El lenguaje expresa intenciones; el code ejecuta acciones. Un modelo que solo maneja lenguaje es como «una persona que sabe pensar y hablar»; solo el Agent que escribe código, invoca herramientas y modifica el entorno empieza a tener manos en el mundo digital.
\end{importantbox}

\begin{knowledgebox}{Términos clave: Affordance}
Affordance proviene de la psicología ecológica y designa las posibilidades de acción que el entorno ofrece a quien actúa. La manija de una puerta ofrece la affordance de «abrir la puerta»; un botón ofrece la affordance de «hacer clic»; para la IA, el código, las API, la línea de comandos y el navegador ofrecen la affordance de «modificar el mundo digital».
\end{knowledgebox}

\begin{knowledgebox}{Niveles de affordance en el mundo digital}
\begin{tabularx}{\textwidth}{p{0.20\textwidth}p{0.34\textwidth}X}
\toprule
Nivel & Ejemplos & Qué puede hacer el Agent \\
\midrule
Nivel de texto & Documentos, páginas web, correos & Leer, resumir, reescribir, buscar. \\
Nivel de código & Python, SQL, shell, API & Calcular, consultar, invocar servicios, modificar el estado del sistema. \\
Nivel de interfaz & Navegador, IDE, sistema operativo & Ejecutar tareas a través de varias herramientas y manejar entornos no estructurados. \\
Nivel organizacional & Colaboración entre varias personas, permisos, aprobaciones & Completar tareas de alto riesgo junto con los procesos humanos. \\
\bottomrule
\end{tabularx}
\end{knowledgebox}

\subsection{Resumen de la sección}

La clave de un sistema de Agent es el diseño conjunto de tareas, entornos, métodos y affordance. El modelo de lenguaje aporta la capacidad cognitiva; el código y las herramientas, la capacidad de acción; el entorno y la reward, la dirección del aprendizaje.
\section{Mecanismos de recompensa: reward intrínseca y Multi-Agent}

Esta sección aborda la reward. Yao Shunyu (姚顺雨) mencionó que una de las direcciones clave del desarrollo de los Agent es que tengan su propia reward y puedan explorar por sí mismos; otra dirección es Multi-Agent, en la que los agentes forman una estructura organizativa entre sí. Aquí la reward no es solo una fórmula de RL: es lo central para que el Agent obtenga retroalimentación efectiva del entorno, vaya más allá de imitar datos humanos y se mejore a sí mismo durante la tarea.

\begin{figure}[H]
\centering
\includegraphics[width=0.94\textwidth]{reward-multi-agent.png}
\caption{Dos direcciones clave para los Agent: tener su propia reward y que múltiples agentes formen una estructura organizativa. Diagrama conceptual propio, elaborado a partir de la conversación de 00:27:19--01:03:00.}
\end{figure}

\begin{knowledgebox}{Lectura de la figura: la reward y el multi-agent son dos vías de expansión distintas}
La reward permite que un solo Agent explore, pruebe, se equivoque y optimice; Multi-Agent permite que varios Agent formen organizaciones más complejas mediante la división del trabajo, la colaboración y la competencia. La primera resuelve la señal de aprendizaje; la segunda, la estructura organizativa.
\end{knowledgebox}

\subsection{Dónde está la dificultad de la reward}

La sección anterior separó la reward y el multi-agent en dos vías de expansión; esta sección examina primero por qué la reward es difícil. En las tareas del mundo real, la reward es difícil de definir. Los juegos de tablero tienen un ganador y un perdedor claros, las tareas de código tienen pruebas y las tareas web pueden tener un estado de finalización; pero los objetivos de muchas tareas reales son de largo plazo, dispersos, subjetivos o no completamente observables. Si la reward está mal diseñada, el Agent puede optimizar el objetivo equivocado o aprovechar los huecos de la evaluación.

\begin{figure}[H]
\centering
\includegraphics[width=0.96\textwidth]{reward-definition-problem.png}
\caption{El problema de definir la recompensa: en las tareas reales, la reward es difícil porque es de largo plazo, dispersa y no completamente observable. Diagrama conceptual propio, elaborado a partir de la conversación de 00:40:00--01:03:00.}
\end{figure}

\begin{warningbox}{El riesgo del reward hacking}
Si la reward solo premia el resultado superficial, el Agent puede aprender a hacer trampa. Por ejemplo, si solo se premia la tasa de clics, puede fabricar titulares sensacionalistas; si solo se premia completar la tarea, puede saltarse los permisos; si solo se premia pasar las pruebas, puede escribir código frágil. Las tareas empresariales y del mundo real requieren en particular auditoría y restricciones de seguridad.
\end{warningbox}

\begin{importantbox}{La fórmula intuitiva de la reward}
\[
\text{learning signal} = r(s_t, a_t, s_{t+1})
\]
Donde \(s_t\) es el estado antes de actuar, \(a_t\) es la acción del Agent, \(s_{t+1}\) es el estado después de actuar y \(r\) es la función de recompensa que evalúa ese cambio. La dificultad de las tareas reales está en que el estado es incompleto, las consecuencias de las acciones llegan con retraso y los criterios de evaluación suelen definirlos en conjunto las personas y las organizaciones.
\end{importantbox}

\subsection{Multi-Agent: de la capacidad individual a la estructura organizativa}

Lo central de Multi-Agent no es «abrir varias instancias del modelo», sino que los Agent formen entre sí división del trabajo, comunicación, colaboración, competencia y una estructura organizativa. Esto se parece a una organización humana y también trae problemas nuevos: quién asigna las tareas, quién verifica los resultados, quién asume los conflictos y cómo evitar las alucinaciones colectivas.

\begin{knowledgebox}{Términos clave: Reward y Multi-Agent}
\begin{tabularx}{\textwidth}{p{0.22\textwidth}p{0.32\textwidth}X}
\toprule
Concepto & Problema que resuelve & Riesgos comunes \\
\midrule
External reward & Evaluación externa que da el entorno o una persona & Costosa, dispersa y propensa al sobreajuste. \\
Intrinsic reward & El Agent genera su propia señal de exploración & Puede desviarse de los objetivos humanos. \\
Multi-Agent & Varios Agent dividen el trabajo y colaboran & Costo de comunicación, conflictos y límites de responsabilidad. \\
Verification & Verificar si la salida del Agent es confiable & El propio verificador también puede equivocarse. \\
\bottomrule
\end{tabularx}
\end{knowledgebox}

La dificultad de la reward y el multi-agent también está en que cambian el «carácter social» del sistema. Cuando un solo Agent se equivoca, el problema puede rastrearse hasta una sola política; cuando varios Agent colaboran, el error puede venir de la asignación de tareas, de malentendidos en la comunicación, de óptimos locales, de defectos del verificador o de incentivos desalineados. Cuanto más se acerca a una organización real, más necesita el sistema de Agent una capa de gobernanza, y no solo un modelo más potente.

\begin{warningbox}{Más agentes no significa automáticamente más capacidad}
Un sistema de múltiples Agent puede aportar división del trabajo y paralelismo, pero también puede traer mayores costos de comunicación, modos de falla más complejos y responsabilidades poco claras. Solo cuando la tarea es descomponible por naturaleza, el mecanismo de verificación es claro y el protocolo de colaboración es estable, es más probable que un sistema multiagente supere a un solo Agent.
\end{warningbox}

\subsection{Resumen de la sección}

La reward determina qué aprende el Agent; Multi-Agent determina cómo se organizan los Agent. La dificultad de la segunda mitad no es que el modelo converse mejor, sino que el sistema pueda explorar y colaborar con seguridad en entornos reales.
\section{Plano de diseño de un sistema de Agent: del concepto de investigación al producto ejecutable}

Las secciones anteriores trataron por separado el lenguaje, el entorno de la tarea, el reward y el multi-agent; esta sección los integra en un plano de sistema que puede diseñarse. Un producto de Agent no suele ser un solo prompt, sino un ciclo: el objetivo del usuario entra al sistema, el sistema convierte ese objetivo en un estado de la tarea, el modelo genera un plan, las herramientas ejecutan acciones, el entorno devuelve observaciones, el verificador evalúa el resultado, la capa de memoria guarda la experiencia y, por último, se pasa a la siguiente ronda.

\begin{importantbox}{Ciclo cerrado mínimo de un sistema de Agent}
\[
\text{Goal} \rightarrow \text{State} \rightarrow \text{Plan} \rightarrow \text{Action} \rightarrow \text{Observation} \rightarrow \text{Reward/Verification} \rightarrow \text{Memory}
\]
Aquí, Goal es el objetivo del usuario, State es el estado actual de la tarea, Plan es el plan que genera el modelo, Action es la ejecución de herramientas o de código, Observation es la retroalimentación del entorno, Reward/Verification es la señal de evaluación y Memory es la experiencia a largo plazo junto con el contexto.
\end{importantbox}

\begin{knowledgebox}{Plano de diseño: qué pregunta cada capa}
\begin{tabularx}{\textwidth}{p{0.20\textwidth}p{0.34\textwidth}X}
\toprule
Capa & Pregunta de diseño & Fallo frecuente \\
\midrule
Capa de objetivo & ¿Qué quiere realmente el usuario? ¿Cuál es la condición de éxito? & El objetivo es ambiguo y el Agent hace algo que parece útil pero no tiene relación con lo pedido. \\
Capa de estado & ¿Qué archivos, herramientas, permisos y restricciones hay en este momento? & Falta contexto y las acciones se basan en un estado erróneo. \\
Capa de planificación & ¿Hace falta dividir la tarea, trabajar en paralelo o consultar al usuario? & El plan es demasiado largo y frágil, y no se ajusta ante una excepción. \\
Capa de acción & ¿Qué acciones pueden ejecutarse con seguridad y cuáles requieren confirmación? & Exceso de permisos, borrados por error, llamadas a la herramienta equivocada. \\
Capa de verificación & ¿Cómo se sabe que la tarea realmente se completó? & Solo se revisa la salida superficial, sin verificar el resultado real. \\
Capa de memoria & ¿Qué experiencia debe conservarse para la próxima vez? & Memoria contaminada, riesgos de privacidad, información desactualizada. \\
\bottomrule
\end{tabularx}
\end{knowledgebox}

\begin{warningbox}{La distancia entre una demo y un producto}
Una demo de Agent puede impresionar con una sola trayectoria exitosa; un producto de Agent tiene que manejar la recuperación ante fallos, los permisos, los registros, la verificación, las interrupciones del usuario, la memoria a largo plazo y el costo. Buena parte de las capacidades de sistema de la segunda mitad reside precisamente en estas capas «aburridas pero necesarias».
\end{warningbox}

\subsection{Resumen de la sección}

Lo esencial del diseño de un sistema de Agent no es envolver el modelo con una UI, sino cerrar el ciclo entre objetivo, estado, plan, acción, observación, verificación y memoria. Solo cuando el ciclo es estable puede el Agent pasar de ser una demostración deslumbrante a un producto confiable.
\section{Frontera de absorción: Interface, Super App y de Chatbot a Agent}

La sección anterior trató los mecanismos internos del Agent; esta sección trata sus fronteras externas. Yao Shunyu (姚顺雨) considera que la mayor oportunidad para las empresas emergentes está en diseñar interfaces distintas. Las capacidades de los modelos pueden dar lugar a formas de interacción beyond ChatGPT y convertirse en una nueva Super App. Aquí, «frontera de absorción» no significa que una empresa vaya a absorberlo todo, sino que cada interfaz de interacción determina en qué tareas puede entrar un sistema inteligente.

\begin{figure}[H]
\centering
\includegraphics[width=0.96\textwidth]{interface-superapp.png}
\caption{Frontera de absorción: la interface determina la oportunidad; distintas formas de interacción pueden dar lugar a distintas Super Apps. Diagrama conceptual propio, elaborado a partir de la conversación entre 00:48:38 y 01:05:01.}
\end{figure}

\begin{knowledgebox}{Lectura de la figura: la interface es una frontera, no solo una UI}
Chatbot, Assistant, Her, interfaces no humanoides, el navegador, el IDE y el sistema operativo pueden convertirse en puntos de entrada a la inteligencia. Cada interface orienta hacia tareas, datos y modelos de negocio distintos, y también moldea la dirección de la investigación.
\end{knowledgebox}

\subsection{Por qué los sistemas Chatbot evolucionan de forma natural hacia Agent}

El sistema Chatbot de una empresa de modelos tiene, por naturaleza, intención del usuario, contexto, necesidad de llamar herramientas y necesidad de estado a largo plazo. Por eso evoluciona de forma natural hacia un sistema Agent: primero recuerda al usuario, luego llama herramientas, después ejecuta tareas y, por último, escribe la retroalimentación de vuelta en el sistema. El Chatbot no es el punto final, sino una de las formas de entrada al Agent.

\begin{figure}[H]
\centering
\includegraphics[width=0.96\textwidth]{chatbot-to-agent.png}
\caption{De Chatbot a Agent: el sistema de chat de una empresa de modelos evoluciona de forma natural hacia un sistema Agent. Diagrama conceptual propio, elaborado a partir de la conversación entre 01:49:28 y 02:05:00.}
\end{figure}

\begin{importantbox}{La transición mínima de Chatbot a Agent}
Basta con que un Chatbot incorpore memory a largo plazo, llamadas a herramientas, estado de la tarea y retroalimentación de resultados para que empiece a convertirse en Agent. La diferencia no está en si la interfaz sigue siendo una ventana de chat, sino en si el sistema puede modificar de manera continua el estado externo.
\end{importantbox}

\subsection{La Super App es un arma de doble filo}

La subsección anterior explicó cómo un Chatbot evoluciona de forma natural hacia un Agent; esta subsección examina la dependencia de trayectoria organizacional que genera un punto de entrada fuerte. Tener una Super App como ChatGPT es una ventaja enorme, porque aporta usuarios, datos, distribución y retroalimentación del producto; pero también moldea la dirección de la investigación y hace que la organización se optimice de forma natural en torno a esa Super App. Si una empresa emergente solo copia una interface existente, difícilmente la superará; si encuentra una interface distinta, puede abrir una nueva frontera de tareas.

\begin{warningbox}{El bloqueo de interfaz influye en la línea de investigación}
Cuando una empresa tiene un punto de entrada fuerte, su investigación tiende a ponerse al servicio de ese punto de entrada. Esto no es negativo, pero la lleva a subestimar las oportunidades de otras interfaces. Históricamente, la sustitución de una plataforma no suele ser una versión mejor de la plataforma anterior, sino una forma de interacción completamente distinta.
\end{warningbox}

\begin{knowledgebox}{Tabla de tipos de Interface}
\begin{tabularx}{\textwidth}{p{0.22\textwidth}p{0.32\textwidth}X}
\toprule
Interface & Tareas típicas & Oportunidad potencial \\
\midrule
Chatbot & Preguntas y respuestas, redacción, explicación, tareas ligeras & Punto de entrada general fuerte, pero propenso a la homogeneización. \\
IDE/Code & Programación, depuración, herramientas de automatización & affordance fuerte; tareas verificables. \\
Browser/OS & Acciones digitales entre páginas web, archivos y aplicaciones & Puede dar lugar a un universal digital agent. \\
Non-human UI & Interacción no humanoide, integrada o en segundo plano & Puede evitar el bloqueo de trayectoria del Chatbot. \\
\bottomrule
\end{tabularx}
\end{knowledgebox}

Por lo tanto, la «frontera de absorción» puede entenderse como la frontera de la interface. Si un modelo solo puede trabajar dentro de una ventana de chat, lo que absorbe son las preguntas y respuestas, la redacción y las tareas ligeras; si entra al IDE, absorbe la producción de código; si entra al navegador y al sistema operativo, absorbe las acciones digitales entre aplicaciones; si entra a los procesos empresariales, absorbe la coordinación repetitiva dentro de la organización. La expansión de la frontera no la logra el modelo por sí solo, sino el modelo, la interfaz, las herramientas y los datos en conjunto.

\subsection{Resumen de la sección}

La frontera de la inteligencia no la determina un solo modelo, sino la interface, las tareas, las herramientas, los hábitos de los usuarios y la trayectoria de la organización en conjunto. Las oportunidades del Agent también surgirán de nuevas formas de interacción.
\section{La humanidad en su conjunto: unipolar y plural a la vez}

La sección anterior analizó cómo la interface determina los límites de la inteligencia; esta sección cierra con la competencia entre plataformas y el panorama de la humanidad en su conjunto. Yao Shunyu (姚顺雨) considera que OpenAI podría convertirse en una empresa tan importante como Google, pero eso no significa que el mundo vaya a quedar monopolizado por un sistema unipolar. Los modelos, la capacidad de cómputo, los datos y las Super App tenderán a concentrarse, y esa es la tendencia unipolar; sin embargo, las distintas interface, las distintas tareas, las distintas estructuras organizacionales y las different bet generan, a su vez, oportunidades plurales.

\begin{figure}[H]
\centering
\includegraphics[width=0.94\textwidth]{single-pole-vs-plural.png}
\caption{Unipolar y plural a la vez: las empresas de modelos importantes serán un eslabón decisivo, pero el mundo no tiene por qué quedar monopolizado por un solo polo. Diagrama conceptual propio, elaborado a partir de la conversación de 01:35:00--02:31:20.}
\end{figure}

\begin{knowledgebox}{Lectura de la figura: lo unipolar y lo plural pueden coexistir}
Los modelos base y los grandes puntos de entrada pueden estar muy concentrados; pero los sistemas inteligentes llegan al mundo por muchas vías, y las diferencias entre tareas, interfaces, herramientas, industrias y organizaciones crean continuamente nuevos puntos de entrada. Por eso, «tener un centro fuerte» y «tener un ecosistema plural» no se contradicen.
\end{knowledgebox}

\begin{knowledgebox}{El problema de cómo hacer allocate de 50,000 millones de dólares}
\begin{tabularx}{\textwidth}{p{0.20\textwidth}p{0.34\textwidth}X}
\toprule
Dirección de la apuesta & Qué podría obtenerse & Riesgo \\
\midrule
Frontier model & Capacidad del modelo y escala de cómputo & Competencia frontal con los gigantes actuales; densidad de capital altísima. \\
Super App & Punto de entrada de usuarios y retorno de datos & Fácil de quedar arrastrada por los puntos de entrada fuertes que ya existen; fuerte dependencia de la trayectoria. \\
Different interface & Nuevas tareas, nuevas interacciones, nuevos ecosistemas & Incertidumbre alta, pero el techo podría ser más alto. \\
Public-good research & Seguridad, evaluación, entornos abiertos y herramientas & Rendimiento comercial lento, pero una contribución más directa a la humanidad. \\
\bottomrule
\end{tabularx}
\end{knowledgebox}

\subsection{Different Bet: superar al líder dominante exige apostar distinto}

Antes se dijo que el mundo tiene a la vez una tendencia a la concentración y oportunidades plurales; esta subsección explica de dónde surgen esas oportunidades plurales. Tanto al inicio como al final de la entrevista aparece una y otra vez la idea de different bet. Yao Shunyu pone un ejemplo: si OpenAI se hubiera dedicado siempre al aprendizaje por refuerzo, difícilmente habría superado a DeepMind; para superar al líder dominante, por lo general hay que apostar por un camino distinto. Las oportunidades de emprendimiento actuales son similares: limitarse a copiar ChatGPT o Claude es muy difícil; si se diseñan Super App distintas, interface distintas o entornos de tareas distintos, todavía hay espacio.

\begin{figure}[H]
\centering
\includegraphics[width=0.96\textwidth]{different-bet.png}
\caption{Different Bet: superar al líder dominante exige apostar distinto, no copiar la línea principal existente. Diagrama conceptual propio, elaborado a partir de la conversación de 00:00:11--00:01:23 y 02:30:20--02:30:37.}
\end{figure}

\begin{importantbox}{Cómo evaluar una oportunidad de emprendimiento}
El verdadero peligro no es que «algo parecido a WeChat (微信) derrote a WeChat», sino que algo distinto derrote a WeChat. En la era de la IA, las nuevas oportunidades también tienen más probabilidades de surgir de nuevas tareas, nuevas interfaces y nuevas organizaciones que de copiar los productos de chat existentes.
\end{importantbox}

\begin{knowledgebox}{Lista de criterios para una Different Bet}
Para saber si una apuesta es realmente different, pueden plantearse cuatro preguntas: ¿cambia el punto de entrada de los usuarios? ¿Cambia el entorno de la tarea? ¿Cambia la retroalimentación y la reward? ¿Cambia la estructura organizacional o el modelo de negocio? Si solo le pone otra apariencia al mismo cuadro de chat, por lo general no es una different bet.
\end{knowledgebox}

\subsection{Resumen de la sección}

«Unipolar y plural a la vez» es la evaluación sobre plataformas más importante de este episodio. Las empresas con modelos fuertes serán nodos decisivos, pero los límites con los que los sistemas inteligentes llegan finalmente al mundo seguirán determinándose por las distintas interface y las different bet.
\section{Síntesis y ampliación}

Esta sección condensa EP115 en cuatro conclusiones. Primera: el lenguaje es una herramienta de generalización, por lo que los modelos de lenguaje resultan naturalmente adecuados como representación de tareas y como interfaz de herramientas para un Agent. Segunda: la investigación sobre Agent no puede fijarse solo en los métodos; también debe atender a las tareas y los entornos. Sin un buen entorno, un método no puede demostrar su valor en el mundo real. Tercera: reward y multi-agent son direcciones importantes de la segunda mitad, pero también traen problemas de seguridad, verificación y organización. Cuarta: la interface determina la frontera por la que la inteligencia entra en el mundo; una Super App es a la vez una ventaja y una dependencia de la trayectoria.

\begin{importantbox}{Situar EP115 en la serie de IA de Zhang Xiaojun (张小珺)}
EP139 trata la historia técnica de los Agent; EP138, el postentrenamiento de Agent y la horizontalidad organizacional; EP116, el Agentic Model para empresas. EP115, en cambio, ofrece un marco de investigación más fundamental: tareas/entornos, métodos simples y generales, reward, multi-agent, code affordance e interface.
\end{importantbox}

\begin{knowledgebox}{Relación con los episodios anteriores y posteriores}
\begin{tabularx}{\textwidth}{p{0.18\textwidth}p{0.34\textwidth}X}
\toprule
Episodio & Foco & Conexión con EP115 \\
\midrule
EP139 & Historia técnica de los Agent y el OpenClaw moment & Ofrece la genealogía histórica de los Agent; EP115 ofrece el marco de investigación. \\
EP138 & Postentrenamiento de Agent, entornos y rollout & Desarrolla por qué la segunda mitad necesita entornos y retroalimentación. \\
EP116 & Agentic Model para empresas y ejecución confiable & Lleva el sistema de Agent de EP115 a los flujos de trabajo ToB. \\
EP118 & Agent OS, VLA y terminales de plataforma & Lleva la interface y los sistemas de Agent al mundo físico y al sistema operativo. \\
\bottomrule
\end{tabularx}
\end{knowledgebox}

\subsection{Takeaways principales}

\begin{enumerate}
\item Un Agent es un sistema capaz de interactuar con un entorno y optimizar un objetivo, no una simple función de chat.
\item Las tareas y los entornos importan tanto como los métodos, e incluso suelen determinar si una investigación tiene sentido en el mundo real.
\item Code es la affordance esencial de la IA en el mundo digital: el puente entre el lenguaje y la acción.
\item Reward es difícil de definir y Multi-Agent es difícil de organizar; ambos son problemas centrales de la segunda mitad.
\item Las nuevas oportunidades surgen de una different bet: interfaces distintas, tareas distintas, organizaciones distintas, y no de copiar al actor dominante.
\end{enumerate}

\subsection{Preguntas abiertas}

Los takeaways anteriores son conclusiones con un grado de certeza relativamente alto; esta sección conserva varias preguntas que siguen abiertas. Son importantes porque la segunda mitad de los Agent no es un camino ya pavimentado: reward, multi-agent, interface y la estructura de las plataformas siguen cambiando con rapidez, y si cambia la respuesta a cualquiera de estas preguntas, puede cambiar la ruta de la siguiente generación de productos e investigación.

\begin{enumerate}
\item ¿Cómo debe diseñarse el reward intrínseco de un Agent para que fomente la exploración sin desviarse de los objetivos humanos?
\item ¿Cómo se reparte la responsabilidad en un sistema Multi-Agent, y quién verifica el resultado final?
\item ¿Las interfaces como Chatbot, IDE, Browser, OS y los robots darán lugar a distintos tipos de Super App?
\item Cuando las empresas de modelos tienen puntos de entrada fuertes, ¿cómo evitan quedar atadas a la interface existente?
\item Si el mundo es a la vez unipolar y plural, ¿en qué tareas y modos de interacción deberían apostar las empresas emergentes?
\end{enumerate}

\begin{warningbox}{Las preguntas abiertas no son un adorno final}
Estas preguntas influyen directamente en la investigación y en los productos: el reward determina la dirección del entrenamiento, el multi-agent determina la forma de organización, la interface determina los puntos de entrada y los datos, y la estructura de las plataformas determina las oportunidades para emprender. Conservar las preguntas es más honesto que ofrecer de antemano una conclusión atractiva pero frágil.
\end{warningbox}

\subsection{Lecturas adicionales}

\begin{itemize}
\item Al contrastarlo con el panorama de Agent de EP139, el marco de investigación de Yao Shunyu (姚顺雨) puede situarse en una historia técnica más larga de los Agent.
\item Al contrastarlo con la entrevista a Luo Fuli (罗福莉) en EP138, se entiende por qué en la etapa de los Agent cobran más importancia el postentrenamiento, los entornos y el rollout.
\item Al contrastarlo con la entrevista a Wu Minghui (吴明辉) en EP116, se observa la versión del Agentic Model para servicios empresariales, con datos privados y ejecución confiable.
\end{itemize}

\end{document}
